\section{Randomized trace estimation}
\label{sec:randomized_trace_estimation}
We have shown that the main quantity of interest in this work is the per-wavenumber mean energy evolution, which, using the cyclic invariance of the trace, can be expressed as
\begin{equation}
    g(k,t)
    =
    \frac{
      \tr
      \left\{
        \mathsfbi{T}\mathsfbi{C}_{00}^{z}
      \right\}
    }{
      \tr
      \left\{
        \mathsfbi{C}_{00}^{z}
      \right\}
    },
    \qquad
    \mathsfbi{T}
    =
    \boldsymbol{\Phi}_{t,k}^{*}\mathsfbi{W}\boldsymbol{\Phi}_{t,k}.
    \label{eq:T_operator_definition}
\end{equation}
This highlights that the required quantity is the trace of the covariance-weighted operator \(\mathsfbi{T}\mathsfbi{C}_{00}^{z}\), normalized by the known initial vertical variance. The major difficulty in computing this quantity is that the propagator \(\boldsymbol{\Phi}_{t,k}\) is a time dependent map.

This distinction is relevant because much of classical non-modal stability theory, as well as the novel perspective of \citet{frame-towne-2024}, is naturally introduced for autonomous linearized systems. In this setting, optimal growth is obtained by maximizing the energy amplification induced by the propagator, which reduces to the matrix exponential and its adjoint when the operator is time independent \citep{schmid-henningson-2001}. For a non-autonomous linear system, however, the propagator is not available in a closed form, even in simplified settings where two of the three spatial directions are homogeneous, and other strategies are required to compute the corresponding amplification \citep{schmid-2007,rapaka-et-al-2008}. Even for autonomous problems, explicitly forming the propagator becomes computationally prohibitive once only one, or none, of the spatial directions remains homogeneous \citep{theofilis-2011}. Consequently, optimal-growth computations are commonly formulated in matrix-free form, where the propagator and its adjoint are applied through time integration of the direct and adjoint equations, rather than by assembling the full matrix operator \citep{luchini-2000,luchini-bottaro-2014}. This idea underlies time-stepper-based computations of transient global optimal disturbances in boundary-layer flows \citep{monokrousos-et-al-2010}, transient growth in stratified shear layers \citep{kaminski-caulfield-taylor-2014}, and spatial optimal-growth analyses of three-dimensional boundary layers \citep{templemann-hanifi-henningson-2010}.

In summary, in optimal-growth analysis, one seeks the largest singular value of the linearized input--output map, or equivalently the largest eigenvalue of the operator obtained by composing the direct propagator with its adjoint. The adjoint-loop procedure is therefore a matrix-free way of applying this composite amplification operator without ever constructing it explicitly. The present problem is closely related, but the quantity of interest is different. Rather than isolating the most amplified initial disturbance, we seek the mean energy growth associated with random initial perturbations. The problem remains a norm-like computation for the linearized input--output map, but instead of the spectral norm, we seek a Frobenius-type measure of the same propagator. The additional complication is that the propagator is generated by a base state that evolves in time and is not time-periodic, so Floquet theory may not be invoked. Once this analogy is evident, the motivation for seeking a loop-adjoint construction is straightforward, in fact,
\begin{equation}
    \mean{
      \left\langle\hat{\mathbf{q}}(t),\hat{\mathbf{q}}(t)\right\rangle_{W}
    }
    =
    \mean{
      \left\langle
        \boldsymbol{\Phi}_{t,k}\hat{\mathbf{q}}_{0,k},
        \boldsymbol{\Phi}_{t,k}\hat{\mathbf{q}}_{0,k}
      \right\rangle_{W}
    }
    =
    \mean{
      \left\langle
        \hat{\mathbf{q}}_{0,k},
        \boldsymbol{\Phi}_{t,k}^{\dagger}\boldsymbol{\Phi}_{t,k}\hat{\mathbf{q}}_{0,k}
      \right\rangle_{W}
    }
    =
    \tr\left\{
      \mathsfbi{T}\mathsfbi{C}_{00}^{z}
    \right\},
    \label{eq:expectation_to_trace}
\end{equation}
where \(\left\langle\mathbf{a},\mathbf{b}\right\rangle_{W}=\mathbf{a}^{*}\mathsfbi{W}\mathbf{b}\) is the \(\mathsfbi{W}\)-weighted inner product, so that \(\mathcal{E}(t)=\left\langle\hat{\mathbf{q}}(t),\hat{\mathbf{q}}(t)\right\rangle_{W}\), and \(\boldsymbol{\Phi}_{t,k}^{\dagger}\) denotes the adjoint of \(\boldsymbol{\Phi}_{t,k}\) with respect to this product, distinct from the conjugate transpose \(\boldsymbol{\Phi}_{t,k}^{*}\).




The loop-adjoint procedure is typically used inside a power or Krylov iteration \citep{barkeley-blackburn-sherwin-2008} to compute optimal growth. Here, instead, its purpose is to build the machinery that allows us to extract the underlying structure of the adjoint-direct operator \(\mathsfbi{T}\) when applied to the covariance matrix, as reported in equation~\eqref{eq:T_operator_definition}. Algorithm~\ref{alg:loop-adjoint} reports the main building block of the method we present in this work; the derivation details for a linear system of the form in equation~\eqref{eq:semidiscrete_system} are given in Appendix~\ref{sec:appendix_main}, \S\ref{sec:appendix_loop_adjoint}.
\begin{algorithm}[t]
  \caption{Matrix-free loop-direct-adjoint (LDA) action of \(\mathsfbi{T}\)}
  \label{alg:loop-adjoint}

  \begin{algorithmic}[1]
    \Require Initial perturbation \(\hat{\mathbf{q}}_{0,k}\), final time \(T\),
             operators \(\mathsfbi{A}(t)\) and \(\mathsfbi{B}(t)\)
    \Ensure Matrix-free action
      \(\LDA(\hat{\mathbf{q}}_{0,k})
      :=
      \mathsfbi{T}\hat{\mathbf{q}}_{0,k}
      =
      \mathsfbi{B}^{\dagger}(0)\boldsymbol{\psi}(0)\)

    \State \(\hat{\mathbf{q}}(0)\gets\hat{\mathbf{q}}_{0,k}\)
      \Comment{Initial condition}

    \State Solve
      \(\mathsfbi{B}(t)\,\partial_t\hat{\mathbf{q}}(t)
      =
      \mathsfbi{A}(t)\hat{\mathbf{q}}(t)\),
      \(0\leq t\leq T\)
      \Comment{Forward solve}

    \State Define \(\boldsymbol{\psi}(T)\) by
      \(\mathsfbi{B}^{\dagger}(T)\boldsymbol{\psi}(T)
      =
      \hat{\mathbf{q}}(T)\)
      \Comment{Terminal adjoint condition}

    \State Solve
      \(\partial_t
      \left[
        \mathsfbi{B}^{\dagger}(t)\boldsymbol{\psi}(t)
      \right]
      =
      -\mathsfbi{A}^{\dagger}(t)\boldsymbol{\psi}(t)\),
      \(T\geq t\geq 0\)
      \Comment{Backward adjoint solve}

    \State \(\LDA(\hat{\mathbf{q}}_{0,k})
      \gets
      \mathsfbi{B}^{\dagger}(0)\boldsymbol{\psi}(0)\)
      \Comment{Matrix-free action}

    \State \Return \(\LDA(\hat{\mathbf{q}}_{0,k})\)
  \end{algorithmic}
\end{algorithm}
The computation of the Frobenius norm of a matrix, or more directly the computation of the trace of a matrix, has been extensively analyzed in the field of randomized numerical linear algebra \citep{haim-sivan-2011}. These methods replace explicit matrix construction by random probing, where one applies the operator to a set of random vectors and combines the resulting quantities to estimate global properties of the operator. Randomized numerical linear algebra has become a standard framework for approximating large matrix quantities using only matrix-vector products, particularly when the operator is too large to form or factorize explicitly \citep{martinsson-tropp-2020}. In fluid mechanics, related randomized ideas have already been used to reduce the cost of large-scale input--output calculations and modal decompositions. For example, \citet{ribeiro-yeh-taira-2020} introduced randomized resolvent analysis, in which random test matrices are used to sketch the resolvent operator and approximate its dominant forcing-response subspace. More recently, \citet{farghadan-martini-towne-2025} combined randomized singular value decomposition with time stepping to compute resolvent modes in fully three-dimensional flows at reduced memory and computational cost. The use of randomization in the present work is related in spirit but different in purpose. Rather than constructing a low-rank approximation in order to recover dominant resolvent or propagator modes, we use random probes to estimate the trace of the covariance-weighted amplification operator. This is achieved by leveraging the Hutchinson trace estimator, which estimates the trace of an implicitly defined matrix from random quadratic forms \citep{hutchinson-1990}. Later work has refined the theory of such estimators and provided bounds on the number of samples required to obtain a prescribed accuracy \citep{khorasani-szekely-ascher-2015}. More recent algorithms, such as Hutch++, further reduce the sampling cost by separating the trace into a low-rank part, which captures the dominant range of the operator, and a residual part, which is estimated stochastically \citep{meyer-et-al-2020}. Extensions from matrices to integral operators have also been proposed recently by \citet{zvonek-horning-townsend-2025}. In this work, building on these algorithms and on the observation that the loop-adjoint technique can be leveraged to build a sketch of our operator, we propose the stochastic stepping trace estimator (SSTE) reported in algorithm~\ref{alg:stochastic_stepping_trace_estimator}.
\begin{algorithm}[t]
  \caption{Stochastic stepping trace estimator for \(\tr\left\{\mathsfbi{T}\mathsfbi{C}_{00}^{z}\right\}\)}
  \label{alg:stochastic_stepping_trace_estimator}

  \begin{algorithmic}[1]
    \Require Number of samples \(m\), loop-direct-adjoint action \(\LDA(\cdot)\),
             factor \(\mathsfbi{L}\) s.t.
             \(\mathsfbi{C}_{00}^{z}=\mathsfbi{L}\mathsfbi{L}^{*}\)
    \Ensure Trace estimate \(\widehat{t}_{\mathrm{H++}}\)

    \State Define kernel
      \(\mathsfbi{K}(\mathbf{g})
      =
      \mathsfbi{L}^{*}\LDA(\mathsfbi{L}\mathbf{g})\)
      \Comment{Effective operator \(\mathsfbi{K}=\mathsfbi{L}^{*}\mathsfbi{T}\mathsfbi{L}\) acting on a vector \(\mathbf{g}\)}

    \State Set \(s \gets m/3\)
      \Comment{Number of probes per stage}

    \For{\(i=1,\ldots,s\)}
      \State Draw \(\mathbf{g}_i \sim \mathcal{N}(\mathbf{0},\mathsfbi{I})\)
        \Comment{Gaussian probe}
      \State \(\mathbf{y}_i \gets \mathsfbi{K}(\mathbf{g}_i)\)
        \Comment{Sketch action}
    \EndFor

    \State Form
      \(\mathsfbi{Y}
      \gets
      [\,\mathbf{y}_1\ \cdots\ \mathbf{y}_s\,]\)
      \Comment{Range sketch}

    \State Compute
      \(\mathsfbi{Y}=\mathsfbi{Q}\mathsfbi{R}\)
      \Comment{Thin QR factorization}

    \State \(t_d \gets 0\)
    \For{\(j=1,\ldots,s\)}
      \State \(t_d \gets t_d + \mathbf{q}_j^{*}\mathsfbi{K}(\mathbf{q}_j)\)
        \Comment{Low-rank trace contribution}
    \EndFor

    \State \(t_r \gets 0\)
    \For{\(i=1,\ldots,s\)}
      \State Draw \(\mathbf{g}_i \sim \mathcal{N}(\mathbf{0},\mathsfbi{I})\)
        \Comment{Residual probe}
      \State \(\mathbf{r}_i \gets (\mathsfbi{I}-\mathsfbi{Q}\mathsfbi{Q}^{*})\mathbf{g}_i\)
        \Comment{Orthogonal projection}
      \State \(t_r \gets t_r + s^{-1}\mathbf{r}_i^{*}\mathsfbi{K}(\mathbf{r}_i)\)
        \Comment{Stochastic residual trace}
    \EndFor

    \State \(\widehat{t}_{\mathrm{H++}} \gets t_d+t_r\)
      \Comment{Hutch++ estimate}

    \State \Return \(\widehat{t}_{\mathrm{H++}}\)
  \end{algorithmic}
\end{algorithm}
Extensive verification against analytical solutions and reference problems is reported in Appendix~\ref{sec:appendix_main}, \S\ref{sec:appendix_verification}.
